% Part III-A, sections that do not depend on the results (included by main_en.tex).

\section{Task: code repair with visible and held-out tests}\label{sec:task}

\paragraph{Problems.} Training problems come from KodCode-Light-RL-10K~\cite{xu2025kodcode}: Python functions with
\texttt{pytest} unit tests. We keep the problems whose test file tests one function, uses no fixtures or third-party
packages, and has at least three test functions. The first test function is shown to the policy (\emph{visible}); the
others are \emph{held out}. 300 problems (seed 0) are held out as the validation set (\texttt{kodcode\_val}). The test
sets are MBPP+ and HumanEval+~\cite{austin2021program,chen2021codex,liu2023evalplus}. For MBPP+ the three original MBPP
asserts are visible. For HumanEval+ the first three asserts of the original HumanEval test are visible. In both, the
EvalPlus harness with its added tests is held out. A problem is dropped if its reference solution fails any of its own
tests under our harness, and we report how many were dropped. Each program runs in its own \texttt{python -I}
subprocess in an empty temporary directory. It has CPU-time, address-space and file-size limits and a wall-clock
timeout, and runs after the destructive functions of \texttt{os}, \texttt{shutil} and \texttt{subprocess} are disabled
(the ``reliability guard'' of OpenAI's \texttt{human-eval}). This contains buggy code from a small model; it is not a
security boundary.

\paragraph{Episodes.} The policy sees the problem, the function signature and the visible tests, and writes a
solution (at most 768 tokens per attempt). The code runs on the visible tests. If one fails, the policy sees the first
failing test and its error and writes another attempt, up to three in all. Every attempt continues the exact token
sequence of the earlier turns. An episode ends at the first attempt that passes every visible test, or after the third,
and its final code then runs on the held-out tests. The prompts are in
\texttt{experiments/part3a\_code/prompts/code\_prompts.json} and were fixed before any run.

\paragraph{Metrics.} \emph{vis pass}: the final code passes every visible test. \emph{all pass}: it passes every
visible and every held-out test. \emph{hack}: it passes the visible tests and fails a held-out one. This counts an
over-fitted or special-cased solution and also a merely incomplete one: the metric cannot tell the two apart, and we
report it as the gap between what the visible tests certify and what is true. \emph{First-attempt all pass} measures
the policy without repair.

\section{Rewards}\label{sec:rewards}

\begin{itemize}
\item \texttt{vis}: 1 if the final code passes every visible test. This is all an agent that holds only the visible
  tests can check, and it is the reward that invites over-fitting to them.
\item \texttt{vis\_kev}: \texttt{vis} times \kev{}'s $p(\text{yes})$ for \emph{genuine}: ``Does the proposed code solve
  the problem for all valid inputs, instead of special-casing the example tests?'' \kev{} sees the problem, the visible
  tests and the final code, never the held-out tests.
\item \texttt{hidden}: 1 if the final code passes every visible and every held-out test. This is an oracle: what
  \texttt{vis} would be with complete tests.
\end{itemize}
Every attempt of an episode is trained with the episode's advantage, and the feedback tokens are context, not trained.
The update is Part II's: group baseline over 8 episodes per problem, batch scale, token-level loss, truncated importance
sampling, one optimizer step per batch.

\section{Offline judge study}\label{sec:code-study}

Before training, the base policy writes 4 single attempts per problem at temperature 1 on \texttt{kodcode\_val}, MBPP+
and HumanEval+. We run them on both test sets, and both judges answer \emph{genuine} and \emph{correct} for each. The
question that matters for a visible-test reward is whether a judge can separate, among the attempts that pass the
visible tests, those that also pass the held-out tests. We report that AUROC, with the AUROC for all-pass over all
attempts. We also build two kinds of code with known status for MBPP+ and HumanEval+. One is the reference solution.
The other is a \emph{hard-coded} lookup that returns the expected value for each visible test's arguments and
\texttt{None} otherwise. It is kept only if it passes the visible tests, so it is a pure test-special-casing program. We
report how often each judge scores the lookup above the reference solution of the same problem.
